% !TeX root = ../main.tex
\section{Threats to Validity}

\textbf{Construct validity.}
Node, edge, classification, and evidence measures capture detectable topology and explicit rules, not component responsibility, architectural rationale, or runtime quality attributes. Exact-line agreement requires one matching evidence item; it does not measure ranking, supporting-location completeness, annotation comprehension, or reviewer effort. Changed-file agreement checks Git scope, not report sufficiency for an architecture decision. The parsed-file fraction measures AST input, not CPU work. D2 treats annotated file-line-detector-edge tuples as signals, so its specificity covers declared hard negatives, not all unannotated syntax.

\textbf{Internal validity.}
The authors curated the analyzer and both datasets. D1 topology intent was established in Phase 1, but patches and exact evidence lines were refined during Phase 2; P3 reuses those known labels, so its agreement is regression, not independent validation of the Git-diff design. D2 followed v0.1 and guided v0.2, creating tuning and confirmation risks. Schema and mutation checks, isolated patches, annotations, frozen hash-bound artifacts, and raw predictions prevent inconsistency and retrospective number changes but do not provide independent ground-truth judgment.

\textbf{External validity.}
The study contains one small synthetic TypeScript/Node.js system, 20 changes, and 40 controlled detector signals. Phase 3 checks working-tree and feature-branch Git states, not longitudinal history, large monorepositories, human approval, or privileged bypass. A 2026-09-11 operational check found all seven ArchSync repositories public and their \texttt{main} branches protected, but those controls are outside the Phase 1--3 study and do not prove human review. The paper repository previously exposed named manuscript source while venue and repository-visibility decisions were unresolved; later redaction cannot establish prior anonymity. Selected \texttt{fetch}, \texttt{pg}, Redis, and AMQP patterns may not generalize to wrappers, dynamic endpoints, reflection, dependency injection, unconventional monorepos, generated code, other languages, infrastructure-as-code, or runtime-only relationships. Contracts are language-neutral, but only the TypeScript analyzer has empirical evidence.

\textbf{Comparative validity.}
No external architecture-conformance or dependency-analysis tool was executed in the reported study. The v0.1--v0.2 comparison uses two ArchSync revisions on a corpus derived from known v0.1 weaknesses, so it measures within-tool regression only. Import-based tools expose different relation sets from ArchSync's HTTP, database, cache, and messaging detectors. A future comparison must freeze common capabilities, versions, configurations, repositories, and ground truth before inspecting outputs. Until then, the paper makes no claim that ArchSync is more accurate, faster, or more useful than an existing tool.

\textbf{Literature-positioning validity.}
The current Related Work synthesis is narrative and may reflect source-selection and interpretation bias. The independently reviewed protocol and screening codebook are frozen at version 1.0.0, but the official four-source search is still in progress and has no frozen screening population; governance does not make the current citation selection systematic. OpenAlex and Semantic Scholar may differ from Scopus and the Web of Science Core Collection in coverage, abstract availability, indexing latency, metadata quality, and query semantics. The study will report source-specific yields and sentinel recall, not subscription-index equivalence. Until all 24 source-query exports, deduplication, dual screening, quality assessment, and synthesis are complete, the cited literature supports scoped positioning only, not coverage completeness or absence of evidence.

\textbf{Conclusion validity.}
D1 has only five positive changed-node items, 12 positive changed-edge items, seven violation rule sets, and 11 evidence-bearing cases. D2 has 20 positive and 20 hard-negative signals. P3 has 20 cold and 20 warm checks over the same D1 changes. Zero-error agreement on these finite, co-developed sets is a feasibility and regression result, not a statistical claim about a population. Latency percentiles are descriptive measurements from one environment; no confidence interval or significance test is used because the samples are neither independent random draws nor a defined workload population.

\textbf{Reliability.}
Dependencies, inputs, paths, JSON, patches, and artifacts are pinned, hashed, or normalized, and every current input is executed twice. Phase 3 requires cache-miss/hit replay and compares incremental and full scans; raw timings remain retained and recomputable. Windows and GitHub Actions reproduced the functional gate, but timing is Windows-only. The v0.1 challenge is frozen and provenance-validated, not re-executed. These controls reduce, but do not eliminate, platform and toolchain threats.
